\section{Introduction}

\subsection{Motivation}

% Frame with plain text
\begin{frame}{Motivation}
  \begin{itemize}
    \item Rise of World Wide Web led to better collaboration between researchers
    \item Place to meet and exchange information
          %researchers in particular from all over the world can collaborate more easily 
    \item This led to the rise of online Communities of Practice (CoP)
          %what Cops are will be discussed in more detail later.
    \item Success of CoP relies on success factors, which are changing over time
    \item It is difficult to measure and evaluate them
          \begin{itemize}
            \item Domain specific success factors
            \item Time consuming
            \item Require data-management experts %which especially small CoPs lack
          \end{itemize}
  \end{itemize}
\end{frame}


\begin{frame}{Motivation}
  %this is why some continuous success modeling systems have been developed
  %however, the problem is that
  \begin{columns}
    \begin{column}[]{0.5\textwidth}

      Current success modeling systems
      \begin{itemize}
        \item Complicated %and require technical know-how
        \item Not optimized for collaboration
        \item Do not take mobility into consideration
      \end{itemize}

    \end{column}
    \begin{column}[]{0.5\textwidth}
      Chat platforms are used for information exchange
      \begin{itemize}
        \item Intuitive
        \item Collaboration
        \item Optimized for mobile devices
      \end{itemize}
      %This shows how chat platforms could adress the aforementionted issues
    \end{column}
  \end{columns}

\end{frame}


\subsection{Thesis Goals}

\begin{frame}{Thesis Goals}
  \begin{itemize}
    \item Design a chatbot for success modeling and visualizations
    \item Interface which communicates between
          \begin{itemize}
            \item End-User
            \item Existing CIS
                  %Especially for visualizations of success factors of the Community
          \end{itemize}
    \item Bot should communicate wit MobSOS (CCA) %framework
    \item Should be deployed with the Social Bot Framework
  \end{itemize}
\end{frame}

%------------------------------------------------

% Start a new section (text is displayed on top of a frame)
\section{State of the Art}
\section{Concept}
\section{Realization}
\section{Evaluation}
\section{Project Plan}
% Frame with items
\begin{frame}{Frame title}
  \begin{itemize}
    \item First level
          \begin{itemize}
            \item Second level
                  \begin{itemize}
                    \item Third level
                  \end{itemize}
          \end{itemize}
  \end{itemize}
\end{frame}

%------------------------------------------------

\begin{frame}
  \frametitle{Blocks of Highlighted Text}
  \begin{block}{Block 1}
    Lorem ipsum dolor sit amet, consectetur adipiscing elit. Integer lectus nisl, ultricies in feugiat rutrum, porttitor sit amet augue. Aliquam ut tortor mauris. Sed volutpat ante purus, quis accumsan dolor.
  \end{block}

  \begin{block}{Block 2}
    Pellentesque sed tellus purus. Class aptent taciti sociosqu ad litora torquent per conubia nostra, per inceptos himenaeos. Vestibulum quis magna at risus dictum tempor eu vitae velit.
  \end{block}

  \begin{block}{Block 3}
    Suspendisse tincidunt sagittis gravida. Curabitur condimentum, enim sed venenatis rutrum, ipsum neque consectetur orci, sed blandit justo nisi ac lacus.
  \end{block}
\end{frame}

%------------------------------------------------

\begin{frame}
  \frametitle{Multiple Columns}
  \begin{columns}[c] % The "c" option specifies centered vertical alignment while the "t" option is used for top vertical alignment

    \column{.45\textwidth} % Left column and width
    \textbf{Heading}
    \begin{enumerate}
      \item Statement
      \item Explanation
      \item Example
    \end{enumerate}

    \column{.5\textwidth} % Right column and width
    Lorem ipsum dolor sit amet, consectetur adipiscing elit. Integer lectus nisl, ultricies in feugiat rutrum, porttitor sit amet augue. Aliquam ut tortor mauris. Sed volutpat ante purus, quis accumsan dolor.

  \end{columns}
\end{frame}

%------------------------------------------------
\section{Second Section}
%------------------------------------------------

\begin{frame}
  \frametitle{Table}
  \begin{table}
    \begin{tabular}{l l l}
      \toprule
      \textbf{Treatments} & \textbf{Response 1} & \textbf{Response 2} \\
      \midrule
      Treatment 1         & 0.0003262           & 0.562               \\
      Treatment 2         & 0.0015681           & 0.910               \\
      Treatment 3         & 0.0009271           & 0.296               \\
      \bottomrule
    \end{tabular}
    \caption{Table caption}
  \end{table}
\end{frame}

%------------------------------------------------

\begin{frame}
  \frametitle{Theorem}
  \begin{theorem}[Mass--energy equivalence]
    $E = mc^2$
  \end{theorem}
\end{frame}

%------------------------------------------------

\begin{frame}[fragile] % Need to use the fragile option when verbatim is used in the slide
  \frametitle{Verbatim}
  \begin{example}[Theorem Slide Code]
    \begin{verbatim}
\begin{frame}
\frametitle{Theorem}
\begin{theorem}[Mass--energy equivalence]
$E = mc^2$
\end{theorem}
\end{frame}\end{verbatim}
  \end{example}
\end{frame}

%------------------------------------------------

\begin{frame}
  \frametitle{Figure}
  Code to include your own image as figure~\ref{fig:acis} from the `Figure' directory.
  \begin{figure}
    \includegraphics[width=0.5\linewidth]{acis}
    \caption{Acis logo}
    \label{fig:acis}
  \end{figure}
\end{frame}

% Frame with text and picture
\begin{frame}{Frame title}
  \begin{columns}[onlytextwidth]\
    % Text on the left
    \begin{column}{.5\textwidth}
      \begin{itemize}
        \item First level
              \begin{itemize}
                \item Second level
                      \begin{itemize}
                        \item Third level\cite{Weng98}
                      \end{itemize}
              \end{itemize}
      \end{itemize}
    \end{column}
    % Picture on the right
    \begin{column}{.5\textwidth}
      \hfill\raisebox{-10cm}[0pt][10cm]{\includegraphics[height=10cm]{figure.jpg}}
    \end{column}
  \end{columns}
\end{frame}


% Frame with large picture and caption
\begin{frame}
  \begin{figure}
    \includegraphics[width=0.7\textwidth]{title_large}
    \caption{Image description}
  \end{figure}
\end{frame}

%------------------------------------------------

% New section (with charts)
% Note: Using newline in a section is not recommended, this is
% just an example
\section{Section title\newline
  Example of a two-line title}

% Frame with horizontal bar chart
\begin{frame}{Add title}
  \begin{center}
    \begin{tikzpicture}
      \begin{axis}[hor_barchart,symbolic y coords={Kategorie 1,Kategorie 2,Kategorie 3, Kategorie 4},%
          width=0.7\textwidth,height=10cm]
        \addplot [draw opacity=0,fill=rwth-50] coordinates {(4.5,Kategorie 4) (3.5,Kategorie 3) (2.5,Kategorie 2) (4.3,Kategorie 1)};
        \addplot [draw opacity=0,fill=rwth] coordinates {(2.8,Kategorie 4) (1.8,Kategorie 3) (4.4,Kategorie 2) (2.4,Kategorie 1)};
        \legend{Datenreihe 1,Datenreihe 2}
      \end{axis}
    \end{tikzpicture}
  \end{center}
\end{frame}

% Frame with vertical bar chart
\begin{frame}{Add title}
  \begin{center}
    \begin{tikzpicture}
      \begin{axis}[ver_barchart,symbolic x coords={Kategorie 1,Kategorie 2,Kategorie 3, Kategorie 4},%
          width=0.7\textwidth,height=10cm]
        \addplot [draw opacity=0,fill=rwth] coordinates {(Kategorie 4,4.5) (Kategorie 3,3.5) (Kategorie 2,2.5) (Kategorie 1,4.3)};
        \addplot [draw opacity=0,fill=rwth-75] coordinates {(Kategorie 4,2.8) (Kategorie 3,1.8) (Kategorie 2,4.4) (Kategorie 1,2.4)};
        \addplot [draw opacity=0,fill=rwth-50] coordinates {(Kategorie 4,3) (Kategorie 3,0.5) (Kategorie 2,5.4) (Kategorie 1,3.4)};
        \addplot [draw opacity=0,fill=rwth-25] coordinates {(Kategorie 4,3.5) (Kategorie 3,1.5) (Kategorie 2,2.2) (Kategorie 1,4.2)};
        \legend{Datenreihe 1,Datenreihe 2,Datenreihe 3,Datenreihe 4}
      \end{axis}
    \end{tikzpicture}
  \end{center}
\end{frame}

% Frame with pie chart
\begin{frame}{Add title}
  \begin{center}
    \hspace{2cm}
    \begin{tikzpicture}
      [
        pie chart,
        slice type={one}{rwth-25},
        slice type={two}{rwth-50},
        slice type={three}{rwth-75},
        slice type={four}{rwth},
        pie values/.style={font={\Large}},
        scale=5
      ]
      \pie{}{29/one,17/two,23.6/three,30.4/four}
      \legend[xshift=1.5cm,yshift=0.8cm]{{Kategorie 1}/one, {Kategorie 2}/two, {Kategorie 3}/three, {Kategorie 4}/four}
    \end{tikzpicture}
  \end{center}
\end{frame}

%------------------------------------------------

% Citation
\begin{frame}[fragile] % Need to use the fragile option when verbatim is used in the slide
  \frametitle{Citation}
  An example of the \verb|\cite| command to cite within the presentation:\\~

  This statement requires citation \cite{Weng98}.
\end{frame}